\documentclass[10pt, aspectratio=169]{beamer}
\usepackage{../beamer_theme_408}
\title{408考研零基础 --- 四科串联总结}
\subtitle{5-4 零基础阶段总结与提高课预告}
\author{}
\date{}
\begin{document}
\frame{\titlepage}

\begin{frame}{本节目标}
    \begin{enumerate}
        \item 回顾零基础课程全部五个模块的知识体系
        \item 理解四门科目之间的交叉关系
        \item 了解提高课的内容安排
        \item 获取备考建议，规划后续复习
    \end{enumerate}
\end{frame}

\begin{frame}{零基础课程知识总图}
    \begin{center}
    \begin{tikzpicture}[node distance=2mm]
        \tikzset{
            mod/.style={rectangle, draw=primary, fill=white!90!primary, rounded corners,
                        text centered, minimum width=55mm, minimum height=7mm, font=\small\bfseries},
            sub/.style={font=\tiny, text=primary}
        }
        \node[mod, fill=red!10] (m0) {模块0 导学};
        \node[sub, below=of m0] {考试概览、学习路线、C语言速成};

        \node[mod, below=of m0, yshift=-8mm, fill=orange!10] (m1) {模块1 数据结构};
        \node[sub, below=of m1] {线性表 $\to$ 栈/队列 $\to$ 树/二叉树 $\to$ 图 $\to$ 查找 $\to$ 排序};

        \node[mod, below=of m1, yshift=-8mm, fill=yellow!10] (m2) {模块2 计算机组成原理};
        \node[sub, below=of m2] {数据表示 $\to$ 运算 $\to$ 存储 $\to$ 指令 $\to$ CPU $\to$ I/O};

        \node[mod, below=of m2, yshift=-8mm, fill=green!10] (m3) {模块3 操作系统};
        \node[sub, below=of m3] {进程管理 $\to$ 内存管理 $\to$ 文件管理 $\to$ I/O管理};

        \node[mod, below=of m3, yshift=-8mm, fill=blue!10] (m4) {模块4 计算机网络};
        \node[sub, below=of m4] {物理层 $\to$ 链路层 $\to$ 网络层 $\to$ 传输层 $\to$ 应用层};
    \end{tikzpicture}
    \end{center}
\end{frame}

\begin{frame}{模块0 --- 导学回顾}
    \begin{block}{考试概览}
        \begin{itemize}
            \item 408计算机学科专业基础综合，满分150分
            \item 数据结构45分、计组45分、操作系统35分、计网25分
        \end{itemize}
    \end{block}
    \begin{block}{学习路线}
        \begin{itemize}
            \item 基础阶段（现在）：理解概念、建立体系
            \item 强化阶段：刷题巩固、专题突破
            \item 冲刺阶段：全真模考、查漏补缺
        \end{itemize}
    \end{block}
    \begin{block}{C语言速成}
        \begin{itemize}
            \item 指针、结构体、动态内存分配是数据结构的编程基础
            \item 理解地址与引用的概念，为计组和OS打下基础
        \end{itemize}
    \end{block}
\end{frame}

\begin{frame}{模块1 --- 数据结构回顾}
    \begin{block}{核心内容}
        \begin{itemize}
            \item \textbf{线性结构}：顺序表、链表、栈、队列 --- 基础数据结构
            \item \textbf{树与二叉树}：遍历、BST、AVL、哈夫曼树 --- 层次关系建模
            \item \textbf{图}：DFS/BFS、最短路径、最小生成树 --- 网状关系建模
            \item \textbf{查找}：二分查找、BST查找、哈希表 --- O(log n)~O(1)
            \item \textbf{排序}：插入/交换/选择/归并/基数 --- 时间与空间复杂度权衡
        \end{itemize}
    \end{block}
    \vspace{0.3em}
    \alert{各科中均用到数据结构：OS的PCB队列、计网的路由表（Trie树）、计组的页表（数组）}
\end{frame}

\begin{frame}{模块2 --- 计算机组成原理回顾}
    \begin{block}{核心内容}
        \begin{itemize}
            \item \textbf{数据表示}：原码/反码/补码、浮点数IEEE754
            \item \textbf{运算}：ALU、加法器、乘法器
            \item \textbf{存储}：Cache、主存、磁盘，存储层次与映射
            \item \textbf{指令系统}：指令格式、寻址方式、CISC vs RISC
            \item \textbf{CPU}：数据通路、控制单元、流水线
            \item \textbf{I/O}：程序查询、中断、DMA三种方式
        \end{itemize}
    \end{block}
    \vspace{0.3em}
    \alert{计组是理解OS底层实现（中断、MMU、DMA）的硬件基础}
\end{frame}

\begin{frame}{模块3 --- 操作系统回顾}
    \begin{block}{核心内容}
        \begin{itemize}
            \item \textbf{进程管理}：PCB、调度、同步互斥（信号量/PV操作）、死锁
            \item \textbf{内存管理}：分页/分段/段页式、虚拟内存、页面置换
            \item \textbf{文件管理}：目录结构、文件分配方式（连续/链接/索引）、磁盘调度
            \item \textbf{I/O管理}：SPOOLing、缓冲、设备驱动
        \end{itemize}
    \end{block}
    \vspace{0.3em}
    \alert{OS是软硬件桥梁：向下依赖计组的CPU/中断/MMU，向上为应用提供抽象接口}
\end{frame}

\begin{frame}{模块4 --- 计算机网络回顾}
    \begin{block}{核心内容}
        \begin{itemize}
            \item \textbf{物理层}：信号编码、传输介质、接口特性
            \item \textbf{数据链路层}：成帧、差错控制（CRC）、流量控制、介质访问控制（CSMA/CD）
            \item \textbf{网络层}：IP协议、子网划分、路由算法（RIP/OSPF/BGP）
            \item \textbf{传输层}：TCP（可靠传输、拥塞控制）vs UDP
            \item \textbf{应用层}：HTTP、DNS、SMTP、FTP
        \end{itemize}
    \end{block}
    \vspace{0.3em}
    \alert{计网是分布式系统的基础，协议栈体现分层设计的工程思想}
\end{frame}

\begin{frame}{四科交叉关系}
    \begin{center}
    \begin{tikzpicture}[node distance=8mm]
        \tikzset{
            subject/.style={ellipse, draw, thick, minimum width=28mm, minimum height=12mm,
                           text centered, font=\small\bfseries},
            label/.style={font=\tiny, text=primary!70!black}
        }
        \node[subject, fill=red!20] (ds) {数据结构};
        \node[subject, right=of ds, fill=green!20] (os) {操作系统};
        \node[subject, below=of ds, fill=yellow!20] (co) {计算机组成};
        \node[subject, below=of os, fill=blue!20] (cn) {计算机网络};

        \draw[<->, thick, primary] (ds) -- (os) node[midway, above, label] {算法实现};
        \draw[<->, thick, primary] (ds) -- (co) node[midway, left, label] {硬件设计};
        \draw[<->, thick, primary] (co) -- (os) node[midway, right, label] {软硬接口};
        \draw[<->, thick, primary] (os) -- (cn) node[midway, below, label] {网络服务};
        \draw[<->, thick, primary] (co) -- (cn) node[pos=0.3, left, label] {传输介质};
        \draw[<->, thick, primary] (ds) -- (cn) node[pos=0.7, right, label] {路由查找};
    \end{tikzpicture}
    \end{center}
    \vspace{0.5em}
    \begin{itemize}
        \item \textbf{数据结构} --- 算法基础，各科实现的核心工具
        \item \textbf{计算机组成原理} --- 硬件基础，理解计算机运行的根本
        \item \textbf{操作系统} --- 软硬件桥梁，管理资源、提供抽象
        \item \textbf{计算机网络} --- 分布式基础，实现多机通信与协作
    \end{itemize}
\end{frame}

\begin{frame}{提高课内容预告}
    \begin{block}{紧扣408考纲逐条精讲}
        \begin{itemize}
            \item 覆盖考纲全部知识点，不留盲区
            \item 重点难点深度剖析
        \end{itemize}
    \end{block}
    \begin{block}{历年真题分类精讲}
        \begin{itemize}
            \item 按知识点分类，掌握命题规律
            \item 逐个题型突破，总结解题技巧
        \end{itemize}
    \end{block}
    \begin{block}{跨科目综合题突破}
        \begin{itemize}
            \item 类似本模块的五科串联思维训练
            \item 针对408越来越灵活的综合大题
        \end{itemize}
    \end{block}
    \begin{block}{模考冲刺与查漏补缺}
        \begin{itemize}
            \item 全真模拟考试，适应考场节奏
            \item 针对性补齐知识短板
        \end{itemize}
    \end{block}
\end{frame}

\begin{frame}{备考建议}
    \begin{enumerate}
        \item \textbf{合理安排时间}：根据自身基础制定复习计划，零基础学生需更多时间投入
        \item \textbf{反复刷真题}：真题是最好的复习资料，近十年真题至少做两遍
        \item \textbf{建立错题本}：记录错题和易混淆概念，定期回顾
        \item \textbf{重视代码题}：数据结构算法题占比不小，手写代码要熟练
        \item \textbf{理解大于背诵}：408重在理解原理，机械记忆难以应对灵活题型
        \item \textbf{保持节奏}：每天坚持学习，避免三天打鱼两天晒网
    \end{enumerate}
\end{frame}

\begin{frame}{本节总结}
    \begin{enumerate}
        \item 零基础课程覆盖了数据结构、计组、OS、计网四大核心科目
        \item 四科相互交叉：数据结构是算法工具，计组是硬件基础，OS是桥梁，计网是分布式基础
        \item 提高课将逐条精讲考纲、真题分类、综合题突破、模考冲刺
        \item 备考建议：合理安排时间、反复刷真题、建立错题本
    \end{enumerate}
    \vspace{1em}
    \begin{center}
        \textcolor{blue}{\textbf{全系列课程完结，恭喜完成408零基础课程！}}
    \end{center}
\end{frame}
\end{document}
